\documentclass[10pt,a4paper]{article}
\usepackage[T1]{fontenc}
\usepackage[utf8]{inputenc}
\usepackage[margin=2cm,landscape]{geometry}
\usepackage{amsmath,amssymb}
\usepackage{array,booktabs,longtable}
\usepackage{xcolor}
\usepackage{enumitem}
\usepackage[hidelinks]{hyperref}
\newcommand{\lean}[1]{\textcolor{blue!50!black}{\texttt{\detokenize{#1}}}}
\setlength{\LTpre}{6pt}\setlength{\LTpost}{6pt}
\renewcommand{\arraystretch}{1.25}
\begin{document}


\begin{center}{\Large\bfseries Theorem 11: upper bound for the cut distance}\\[4pt]

{\small Correspondence between the paper and the Lean~4/Mathlib formalisation}\end{center}

\noindent\textbf{Paper.} J.~Breitung, T.~Gervens, M.~Grohe. \emph{Graph Similarity Meets Optimal Transport} (file \texttt{Graph\_Similarity\_Meets\_Optimal\_Transport.pdf} in the repository). Statement, equation, page and line numbers refer to this PDF; \texttt{docs/NUMBERING.md} lists the Lean files and declarations of all numbered statements.\par\smallskip

The folder \texttt{GsimOT/Theorem11\_CutUpperBound} formalises Theorem~11 (p.~10), the cut part of Theorem~2, and Appendix~E.2: the multiplicative rounding via thin matchings (E.2.1: Lemmas~36, 37, Theorem~38), the additive rounding (E.2.2: Lemmas~40--45, 47, 48, Corollary~46, Theorem~39), the Talagrand contraction lemma (Lemma~43), and the final proof (E.2.3).\par\paragraph{Checking.} Run \texttt{bash build.sh} in the repository root (see \texttt{README.md}). The script builds the library and runs \texttt{scripts/CheckAxioms.lean}, which must report only the standard axioms \texttt{propext}, \texttt{Classical.choice}, \texttt{Quot.sound}: no \texttt{sorry} and no additional axioms.

\paragraph{Main statement in Lean.} \lean{GsimOT.theorem11_of_thin}: assuming Theorem 35 [HOG26] (\lean{ThinPermutationTheorem}),
there is a universal $C>0$ such that for all graphs $G,H$ of the same order, with $f=\delta^{\mathrm{OT}}_\square(G,H)$,
\[ f\le\delta_\square(G,H)\le C\,f\ln(e/f)\ln\bigl(e\ln(e/f)\bigr). \]
\lean{GsimOT.theorem11_power_lt_one_of_thin}: for every $\alpha<1$ there is $C_\alpha>0$ with
$\delta_\square(G,H)\le C_\alpha\,\delta^{\mathrm{OT}}_\square(G,H)^\alpha$; \lean{GsimOT.theorem2_cut_of_thin}: the same with
$\alpha=1-\varepsilon$ (Theorem~2). The \texttt{\_of\_rackeShah} versions assume only [RS14].

\section*{Correspondence}

\begin{longtable}{>{\raggedright\arraybackslash}p{3.2cm} >{\centering\arraybackslash}p{1.1cm} >{\raggedright\arraybackslash}p{8.4cm} >{\raggedright\arraybackslash}p{10.2cm}}
\toprule
\textbf{Paper} & \textbf{Page} & \textbf{Statement} & \textbf{Lean declaration(s) and location} \\
\midrule\endhead
\bottomrule\endfoot
{}Theorem~11 & 10 & $\delta^{\mathrm{OT}}_\square\le\delta_\square\le C\,\delta^{\mathrm{OT}}_\square\ln\frac{e}{\delta^{\mathrm{OT}}_\square}\ln\bigl(e\ln\frac{e}{\delta^{\mathrm{OT}}_\square}\bigr)$ (assuming Thm.~35) & \lean{GsimOT.theorem11_of_thin}\newline{\scriptsize\texttt{Theorem11\_CutUpperBound/Theorem11.lean:593}} \newline \lean{GsimOT.theorem11}\newline{\scriptsize\texttt{Theorem11\_CutUpperBound/Theorem11.lean:517}} \newline \lean{GsimOT.theorem11_of_rackeShah}\newline{\scriptsize\texttt{Theorem11\_CutUpperBound/FromThinMatching.lean:28}} \\ \midrule
{}Theorem~11 (power bound), E.2.3 & 10, 44 & $\delta_\square\le C_\alpha(\delta^{\mathrm{OT}}_\square)^\alpha$ for every $\alpha<1$ (also the $0<\alpha<1$ versions) & \lean{GsimOT.theorem11_power_lt_one_of_thin}\newline{\scriptsize\texttt{Theorem11\_CutUpperBound/Theorem11.lean:603}} \newline \lean{GsimOT.theorem11_power_lt_one}\newline{\scriptsize\texttt{Theorem11\_CutUpperBound/Theorem11.lean:550}} \newline \lean{GsimOT.theorem11_power_lt_one_of_rackeShah}\newline{\scriptsize\texttt{Theorem11\_CutUpperBound/FromThinMatching.lean:47}} \newline \lean{GsimOT.theorem11_power_of_thin}\newline{\scriptsize\texttt{Theorem11\_CutUpperBound/Theorem11.lean:612}} \newline \lean{GsimOT.theorem11_power}\newline{\scriptsize\texttt{Theorem11\_CutUpperBound/Theorem11.lean:579}} \\ \midrule
{}Theorem~2 (cut part) & 3 & $\delta_\square\le C_\varepsilon(\delta^{\mathrm{OT}}_\square)^{1-\varepsilon}$ for every $\varepsilon>0$ & \lean{GsimOT.theorem2_cut_of_thin}\newline{\scriptsize\texttt{Theorem11\_CutUpperBound/Theorem11.lean:622}} \newline \lean{GsimOT.theorem2_cut_of_rackeShah}\newline{\scriptsize\texttt{Theorem11\_CutUpperBound/FromThinMatching.lean:56}} \\ \midrule
{}Theorem~35 [HOG26] & 32 & Thin permutations exist with $\alpha\le C_{\mathrm{thin}}\ln(2n)\ln\ln(4n)$ --- \textbf{explicit hypothesis} & \lean{GsimOT.ThinPermutationTheorem}\newline{\scriptsize\texttt{Theorem11\_CutUpperBound/Definitions.lean:40}} \newline \lean{GsimOT.IsThin}\newline{\scriptsize\texttt{Theorem11\_CutUpperBound/Definitions.lean:30}} \\ \midrule
{}Theorem~35 from [HOG26, Thm.~1.4] & 32 & Derived from the \texttt{ThinMatching} library, assuming only [RS14] & \lean{GsimOT.thinPermutationTheorem_of_rackeShah}\newline{\scriptsize\texttt{Theorem11\_CutUpperBound/FromThinMatching.lean:18}} \\ \midrule
{}Theorem~39 & 36 & $\zeta_P\le\zeta_Q+36\sqrt{d_n}\,n^{-1/6}$, $d_n=\lceil2\log n\rceil$ ($n\ge2$; $A,B$ symmetric with entries in $[0,1]$) & \lean{GsimOT.theorem39}\newline{\scriptsize\texttt{Theorem11\_CutUpperBound/Theorem39\_AdditiveRounding.lean:267}} \newline \lean{GsimOT.dn}\newline{\scriptsize\texttt{Theorem11\_CutUpperBound/Definitions.lean:76}} \\ \midrule
{}Lemma~40 & 36 & Lifted kernels: $\|K\|_{\mathcal L^2(n)}\le t\Rightarrow\|K^{(1)}\|_{\square,\mu_Q},\|K^{(2)}\|_{\square,\mu_Q}\le t$ & \lean{GsimOT.lemma40_fst}\newline{\scriptsize\texttt{Theorem11\_CutUpperBound/Corollary46\_SpectralTruncation.lean:224}} \newline \lean{GsimOT.lemma40_snd}\newline{\scriptsize\texttt{Theorem11\_CutUpperBound/Corollary46\_SpectralTruncation.lean:258}} \newline \lean{GsimOT.OpBound}\newline{\scriptsize\texttt{Theorem11\_CutUpperBound/Corollary46\_SpectralTruncation.lean:75}} \newline \lean{GsimOT.weight_sq_le}\newline{\scriptsize\texttt{Theorem11\_CutUpperBound/Corollary46\_SpectralTruncation.lean:202}} \\ \midrule
{}Eqs.~(21), (22) & 32 & $\zeta_Q$ (\lean{epsR}); bilinear forms with $[0,1]$-valued functions are bounded by the cut norm; testing on sets & \lean{GsimOT.epsR}\newline{\scriptsize\texttt{Common/CutNorm.lean:105}} \newline \lean{GsimOT.abs_bilin_le_cutNorm}\newline{\scriptsize\texttt{Theorem11\_CutUpperBound/Corollary46\_SpectralTruncation.lean:178}} \newline \lean{GsimOT.abs_cut_le_epsR}\newline{\scriptsize\texttt{Theorem11\_CutUpperBound/Theorem38\_ThinRounding.lean:66}} \\ \midrule
{}Lemma~41 & 37 & Close support functions $\Rightarrow$ close bilinear maxima (paper's proof: $m_T(Y,Y')$ via $h_Y(Tu')$ and $h_{Y'}(T^*u)$) & \lean{GsimOT.lemma41}\newline{\scriptsize\texttt{Theorem11\_CutUpperBound/Corollary46\_SpectralTruncation.lean:588}} \newline \lean{GsimOT.lemma41_step}\newline{\scriptsize\texttt{Theorem11\_CutUpperBound/Corollary46\_SpectralTruncation.lean:532}} \newline \lean{GsimOT.abs_inner_le_bilMax_add}\newline{\scriptsize\texttt{Theorem11\_CutUpperBound/Corollary46\_SpectralTruncation.lean:568}} \newline \lean{GsimOT.abs_supportFn_sub_le}\newline{\scriptsize\texttt{Theorem11\_CutUpperBound/Corollary46\_SpectralTruncation.lean:495}} \newline \lean{GsimOT.supportFn}\newline{\scriptsize\texttt{Theorem11\_CutUpperBound/Corollary46\_SpectralTruncation.lean:461}} \newline \lean{GsimOT.bilMax}\newline{\scriptsize\texttt{Theorem11\_CutUpperBound/Corollary46\_SpectralTruncation.lean:464}} \\ \midrule
{}Lemma~42 & 37 & Symmetrisation (independent $X_a\in\{0,1\}$ with $\mathbb E X_a=\theta_a$) & \lean{GsimOT.Lemma47.lemma42}\newline{\scriptsize\texttt{Theorem11\_CutUpperBound/Lemma47\_MergingPermutations.lean:394}} \newline \lean{GsimOT.Lemma47.bern}\newline{\scriptsize\texttt{Theorem11\_CutUpperBound/Lemma47\_MergingPermutations.lean:321}} \\ \midrule
{}Lemma~43 [LT91, Thm.~4.12] & 38 & Talagrand contraction lemma (cited in the paper; \textbf{proved here}) & \lean{GsimOT.ContractionPrinciple}\newline{\scriptsize\texttt{Theorem11\_CutUpperBound/Definitions.lean:54}} \newline \lean{GsimOT.contractionPrinciple}\newline{\scriptsize\texttt{Theorem11\_CutUpperBound/ContractionPrinciple.lean:173}} \newline \lean{GsimOT.ContractionAux.two_point}\newline{\scriptsize\texttt{Theorem11\_CutUpperBound/ContractionPrinciple.lean:44}} \newline \lean{GsimOT.ContractionAux.comparison}\newline{\scriptsize\texttt{Theorem11\_CutUpperBound/ContractionPrinciple.lean:144}} \\ \midrule
{}Lemma~44, Eq.~(29) & 38--39 & Block Rademacher estimate $\le2M(\sum_j\ell_j^2)^{1/2}$ (random representatives $I_j\in B_j$, Jensen, Lemma~43) & \lean{GsimOT.Lemma47.lemma44}\newline{\scriptsize\texttt{Theorem11\_CutUpperBound/Lemma47\_MergingPermutations.lean:208}} \newline \lean{GsimOT.Lemma47.sum_piFinset_comp}\newline{\scriptsize\texttt{Theorem11\_CutUpperBound/Lemma47\_MergingPermutations.lean:187}} \newline \lean{GsimOT.Lemma47.radAvg_norm_le}\newline{\scriptsize\texttt{Theorem11\_CutUpperBound/Lemma47\_MergingPermutations.lean:85}} \\ \midrule
{}Proof of Lemma~45 & 39--40 & Spectral theorem, Parseval, truncations $A_t,B_t$ with $\|A-A_t\|_{\mathcal L^2(n)},\|B-B_t\|_{\mathcal L^2(n)}\le t$ & \lean{GsimOT.spectral_entry}\newline{\scriptsize\texttt{Theorem11\_CutUpperBound/Corollary46\_SpectralTruncation.lean:48}} \newline \lean{GsimOT.parseval_row}\newline{\scriptsize\texttt{Theorem11\_CutUpperBound/Corollary46\_SpectralTruncation.lean:57}} \newline \lean{GsimOT.specTrunc}\newline{\scriptsize\texttt{Theorem11\_CutUpperBound/Corollary46\_SpectralTruncation.lean:68}} \newline \lean{GsimOT.opBound_sub_specTrunc}\newline{\scriptsize\texttt{Theorem11\_CutUpperBound/Corollary46\_SpectralTruncation.lean:101}} \\ \midrule
{}Eq.~(31) & 39 & $|\zeta_Q(A,B)-\zeta_Q(A_t,B_t)|\le2t$ & \lean{GsimOT.eq31}\newline{\scriptsize\texttt{Theorem11\_CutUpperBound/Corollary46\_SpectralTruncation.lean:296}} \newline \lean{GsimOT.cutNorm_le_of_eq_add3}\newline{\scriptsize\texttt{Theorem11\_CutUpperBound/Corollary46\_SpectralTruncation.lean:139}} \\ \midrule
{}Proof of Lemma~45 & 39--40 & Vectors $z_{ij}$ ($\|z_{ij}\|\le\sqrt2$), operator $T_t$ ($\|T_t\|_{2\to2}\le1/t$), $\langle z_{ij},T_tz_{k\ell}\rangle=(A_t)_{ik}-(B_t)_{j\ell}$ & \lean{GsimOT.specVec}\newline{\scriptsize\texttt{Theorem11\_CutUpperBound/Corollary46\_SpectralTruncation.lean:700}} \newline \lean{GsimOT.specDiag}\newline{\scriptsize\texttt{Theorem11\_CutUpperBound/Corollary46\_SpectralTruncation.lean:708}} \newline \lean{GsimOT.diagOp}\newline{\scriptsize\texttt{Theorem11\_CutUpperBound/Corollary46\_SpectralTruncation.lean:788}} \newline \lean{GsimOT.norm_specVec_le}\newline{\scriptsize\texttt{Theorem11\_CutUpperBound/Corollary46\_SpectralTruncation.lean:754}} \newline \lean{GsimOT.abs_specDiag_le}\newline{\scriptsize\texttt{Theorem11\_CutUpperBound/Corollary46\_SpectralTruncation.lean:731}} \newline \lean{GsimOT.norm_diagOp_le}\newline{\scriptsize\texttt{Theorem11\_CutUpperBound/Corollary46\_SpectralTruncation.lean:808}} \newline \lean{GsimOT.diagForm_specVec}\newline{\scriptsize\texttt{Theorem11\_CutUpperBound/Corollary46\_SpectralTruncation.lean:713}} \newline \lean{GsimOT.inner_diagOp}\newline{\scriptsize\texttt{Theorem11\_CutUpperBound/Corollary46\_SpectralTruncation.lean:799}} \\ \midrule
{}Lemma~45, Eq.~(30) & 39 & Spectral approximation: $|\zeta_Q-\max_{S,T}|\ldots\langle z_{ij},T_tz_{k\ell}\rangle\ldots||\le2t$ & \lean{GsimOT.lemma45}\newline{\scriptsize\texttt{Theorem11\_CutUpperBound/Corollary46\_SpectralTruncation.lean:826}} \newline \lean{GsimOT.zonKernel}\newline{\scriptsize\texttt{Theorem11\_CutUpperBound/Corollary46\_SpectralTruncation.lean:627}} \\ \midrule
{}Eqs.~(32)--(34) & 40 & Zonotope $Z_Q$, its support function $h_Q$ (Eq.~(33)), pseudometric $\rho$ & \lean{GsimOT.zonLin}\newline{\scriptsize\texttt{Theorem11\_CutUpperBound/Corollary46\_SpectralTruncation.lean:336}} \newline \lean{GsimOT.cube}\newline{\scriptsize\texttt{Theorem11\_CutUpperBound/Corollary46\_SpectralTruncation.lean:346}} \newline \lean{GsimOT.hSupp}\newline{\scriptsize\texttt{Theorem11\_CutUpperBound/Definitions.lean:66}} \newline \lean{GsimOT.supportFn_zonotope}\newline{\scriptsize\texttt{Theorem11\_CutUpperBound/Corollary46\_SpectralTruncation.lean:616}} \newline \lean{GsimOT.rhoSupp}\newline{\scriptsize\texttt{Theorem11\_CutUpperBound/Definitions.lean:70}} \newline \lean{GsimOT.norm_zonLin_le}\newline{\scriptsize\texttt{Theorem11\_CutUpperBound/Corollary46\_SpectralTruncation.lean:396}} \newline \lean{GsimOT.isCompact_zonotope}\newline{\scriptsize\texttt{Theorem11\_CutUpperBound/Corollary46\_SpectralTruncation.lean:355}} \\ \midrule
{}Eq.~(35) & 40 & $\rho$ is a pseudometric, convex in both arguments & \lean{GsimOT.rhoSupp_triangle}\newline{\scriptsize\texttt{Theorem11\_CutUpperBound/Theorem39\_AdditiveRounding.lean:61}} \newline \lean{GsimOT.rhoSupp_convex}\newline{\scriptsize\texttt{Theorem11\_CutUpperBound/Theorem39\_AdditiveRounding.lean:77}} \newline \lean{GsimOT.hSupp_lin}\newline{\scriptsize\texttt{Theorem11\_CutUpperBound/Theorem39\_AdditiveRounding.lean:70}} \\ \midrule
{}Eqs.~(36), (37) & 40--41 & $\max_{u,u'\in Z_Q}|\langle u,T_tu'\rangle|$ is the maximum of Eq.~(30) (attained at indicators); $|\zeta_Q-\max_{Z_Q}|\le2t$ & \lean{GsimOT.eq36}\newline{\scriptsize\texttt{Theorem11\_CutUpperBound/Corollary46\_SpectralTruncation.lean:646}} \newline \lean{GsimOT.eq36_expand}\newline{\scriptsize\texttt{Theorem11\_CutUpperBound/Corollary46\_SpectralTruncation.lean:634}} \newline \lean{GsimOT.eq37}\newline{\scriptsize\texttt{Theorem11\_CutUpperBound/Corollary46\_SpectralTruncation.lean:860}} \\ \midrule
{}Corollary~46 & 41 & $\zeta_P\le\zeta_Q+4t+\frac{2\sqrt2}{t}\rho(P,Q)$ for all $P,Q\in\mathrm{DS}(n)$ & \lean{GsimOT.corollary46}\newline{\scriptsize\texttt{Theorem11\_CutUpperBound/Corollary46\_SpectralTruncation.lean:870}} \\ \midrule
{}Lemma~47 & 41 & Merging two permutations: $\rho(P,(1-\theta)P_0+\theta P_1)\le20Mn^{-1/3}$ (the proof gives $18M$) & \lean{GsimOT.lemma47}\newline{\scriptsize\texttt{Theorem11\_CutUpperBound/Lemma47\_MergingPermutations.lean:961}} \newline \lean{GsimOT.lemma47_sharp}\newline{\scriptsize\texttt{Theorem11\_CutUpperBound/Lemma47\_MergingPermutations.lean:810}} \\ \midrule
{}Proof of Lemma~47 & 42 & Blocks along the cycles of $\tau$ with properties (i), (ii) & \lean{GsimOT.Lemma47.exists_sep_cover}\newline{\scriptsize\texttt{Theorem11\_CutUpperBound/Lemma47\_MergingPermutations.lean:508}} \newline \lean{GsimOT.Lemma47.exists_blocks}\newline{\scriptsize\texttt{Theorem11\_CutUpperBound/Lemma47\_MergingPermutations.lean:547}} \\ \midrule
{}Proof of Lemma~47, Eq.~(38) & 42--43 & A good choice $X$ ($\Delta(T_X)\le16Mn^{-1/3}$); repairing $T_X$ into a bijection & \lean{GsimOT.Lemma47.exists_good_choice}\newline{\scriptsize\texttt{Theorem11\_CutUpperBound/Lemma47\_MergingPermutations.lean:680}} \newline \lean{GsimOT.Lemma47.exists_equiv_eq_off}\newline{\scriptsize\texttt{Theorem11\_CutUpperBound/Lemma47\_MergingPermutations.lean:644}} \\ \midrule
{}Proof of Lemma~48, Eq.~(3) & 43 & Birkhoff + Carath\'eodory: $Q=\sum_{s\le N}\lambda_sP_s$, $N\le n^2$ & \lean{GsimOT.exists_perm_decomp}\newline{\scriptsize\texttt{Theorem11\_CutUpperBound/Theorem39\_AdditiveRounding.lean:183}} \\ \midrule
{}Lemma~48 & 43 & Binary-tree merging: $\rho(P,Q)\le20Md_nn^{-1/3}$ & \lean{GsimOT.lemma48}\newline{\scriptsize\texttt{Theorem11\_CutUpperBound/Theorem39\_AdditiveRounding.lean:243}} \newline \lean{GsimOT.merge_aux}\newline{\scriptsize\texttt{Theorem11\_CutUpperBound/Theorem39\_AdditiveRounding.lean:127}} \newline \lean{GsimOT.combo}\newline{\scriptsize\texttt{Theorem11\_CutUpperBound/Theorem39\_AdditiveRounding.lean:103}} \\ \midrule
{}Proof of Theorem~39 & 43 & $t=\sqrt{20d_n}\,n^{-1/6}$: $4t+80d_nn^{-1/3}/t=8t<36\sqrt{d_n}\,n^{-1/6}$ & \lean{GsimOT.theorem39}\newline{\scriptsize\texttt{Theorem11\_CutUpperBound/Theorem39\_AdditiveRounding.lean:267}} \\ \midrule
{}Eq.~(23) & 32 & $\alpha$-thin permutations: thinness for all cuts $(X\cup Y,\overline X\cup\overline Y)$ & \lean{GsimOT.IsThin}\newline{\scriptsize\texttt{Theorem11\_CutUpperBound/Definitions.lean:30}} \\ \midrule
{}Lemma~36 & 33 & Thinness $\Rightarrow\sum_i|x_i-y_{\pi(i)}|\le\alpha\sum_{ij}Q_{ij}|x_i-y_j|$ (level sets, integral) & \lean{GsimOT.lemma36}\newline{\scriptsize\texttt{Theorem11\_CutUpperBound/Theorem38\_ThinRounding.lean:107}} \newline \lean{GsimOT.levelInd}\newline{\scriptsize\texttt{Theorem11\_CutUpperBound/Theorem38\_ThinRounding.lean:85}} \newline \lean{GsimOT.levelInd_apply}\newline{\scriptsize\texttt{Theorem11\_CutUpperBound/Theorem38\_ThinRounding.lean:87}} \newline \lean{GsimOT.integral_levelInd}\newline{\scriptsize\texttt{Theorem11\_CutUpperBound/Theorem38\_ThinRounding.lean:99}} \\ \midrule
{}Lemma~37, Eq.~(24) & 33--34 & Local agreement (for $y=\mathbf1_T$, the case used) & \lean{GsimOT.lemma37}\newline{\scriptsize\texttt{Theorem11\_CutUpperBound/Theorem38\_ThinRounding.lean:139}} \\ \midrule
{}Eq.~(25) & 34 & $\|RK\|_\square\le\|K\|_\square$ for $R\in\mathrm{DS}(n)$ (one-dimensional form) & \lean{GsimOT.abs_sum_mul_le_of_cut}\newline{\scriptsize\texttt{Theorem11\_CutUpperBound/Theorem38\_ThinRounding.lean:30}} \\ \midrule
{}Theorem~38 & 34 & There is $P$ with $\zeta_P\le(1+4\alpha)\zeta_Q$, $\alpha=C_{\mathrm{thin}}\ln(2n)\ln\ln(4n)$ & \lean{GsimOT.theorem38}\newline{\scriptsize\texttt{Theorem11\_CutUpperBound/Theorem38\_ThinRounding.lean:404}} \\ \midrule
{}Proof of Thm.~38, Eq.~(26) & 34 & $\zeta_P\le(1+4\alpha)\zeta_Q$ for every $\alpha$-thin $P$; $\zeta_P=n^{-2}\|A-PBP^\top\|_\square$ & \lean{GsimOT.eq26}\newline{\scriptsize\texttt{Theorem11\_CutUpperBound/Theorem38\_ThinRounding.lean:373}} \newline \lean{GsimOT.epsR_permMat_le}\newline{\scriptsize\texttt{Theorem11\_CutUpperBound/Theorem38\_ThinRounding.lean:232}} \newline \lean{GsimOT.sum_permMat_mul_atoms}\newline{\scriptsize\texttt{Theorem11\_CutUpperBound/Theorem38\_ThinRounding.lean:218}} \\ \midrule
{}Proof of Thm.~38, Eq.~(27) & 35 & $n^{-2}\|Q^\top AQ-B\|_\square\le\zeta_Q$ & \lean{GsimOT.eq27_left}\newline{\scriptsize\texttt{Theorem11\_CutUpperBound/Theorem38\_ThinRounding.lean:257}} \\ \midrule
{}Proof of Thm.~38 & 35 & $n^{-2}\|EA\|_\square\le2\alpha\zeta_Q$ (Lemmas~36, 37, mean zero) & \lean{GsimOT.normEA_le}\newline{\scriptsize\texttt{Theorem11\_CutUpperBound/Theorem38\_ThinRounding.lean:285}} \newline \lean{GsimOT.abs_sum_le_half_of_sum_eq_zero}\newline{\scriptsize\texttt{Theorem11\_CutUpperBound/Theorem38\_ThinRounding.lean:54}} \\ \midrule
{}Proof of Thm.~38, Eqs.~(28), (25) & 35 & $A-RAR^\top=EA+R(EA)^\top$; $\|R(EA)^\top\|_\square\le\|EA\|_\square$ & \lean{GsimOT.eq28_cut}\newline{\scriptsize\texttt{Theorem11\_CutUpperBound/Theorem38\_ThinRounding.lean:340}} \newline \lean{GsimOT.cut_R_EAt_le}\newline{\scriptsize\texttt{Theorem11\_CutUpperBound/Theorem38\_ThinRounding.lean:311}} \\ \midrule
{}Proof of Thm.~38, Eq.~(27) & 35 & $n^{-2}\|A-RAR^\top\|_\square\le4\alpha\zeta_Q$ & \lean{GsimOT.eq27_right}\newline{\scriptsize\texttt{Theorem11\_CutUpperBound/Theorem38\_ThinRounding.lean:355}} \\ \midrule
{}Section~2; proof of Thm.~11 & 5, 44 & $C_\square(P_\pi)=\|A^\pi_G-A_H\|_\square$; lower bound $\delta^{\mathrm{OT}}_\square\le\delta_\square$ & \lean{GsimOT.costCut_permMat}\newline{\scriptsize\texttt{Theorem11\_CutUpperBound/Theorem11.lean:348}} \newline \lean{GsimOT.deltaOTCut_le_deltaCut}\newline{\scriptsize\texttt{Theorem11\_CutUpperBound/Theorem11.lean:455}} \\ \midrule
{}Proof of Thm.~11 & 44 & An optimal $Q$ exists (compactness of $\mathrm{DS}$, continuity of $C_\square$) & \lean{GsimOT.isCompact_DS}\newline{\scriptsize\texttt{Theorem11\_CutUpperBound/Theorem11.lean:384}} \newline \lean{GsimOT.continuous_costCut}\newline{\scriptsize\texttt{Theorem11\_CutUpperBound/Theorem11.lean:404}} \newline \lean{GsimOT.exists_costCut_eq_dOTCut}\newline{\scriptsize\texttt{Theorem11\_CutUpperBound/Theorem11.lean:416}} \\ \midrule
{}Eq.~(39) & 44 & $\delta_\square\le C_1\ln(2n)\ln\ln(4n)\,\varepsilon$, $C_1=1+4C_{\mathrm{thin}}$, using $\ln4\cdot\ln\ln8>1$ & \lean{GsimOT.eq39}\newline{\scriptsize\texttt{Theorem11\_CutUpperBound/Theorem11.lean:491}} \newline \lean{GsimOT.one_le_log_mul_log_log}\newline{\scriptsize\texttt{Theorem11\_CutUpperBound/Theorem11.lean:281}} \newline \lean{GsimOT.one_lt_log_four_mul_log_log_eight}\newline{\scriptsize\texttt{Theorem11\_CutUpperBound/Theorem11.lean:260}} \\ \midrule
{}Proof of Thm.~11 & 44 & Case analysis ($\varepsilon=0$, $\varepsilon\ge\frac12$, $n\le\varepsilon^{-12}$, $n>\varepsilon^{-12}$), $C=\max\{2,39C_1,113\}$ & \lean{GsimOT.main_numeric}\newline{\scriptsize\texttt{Theorem11\_CutUpperBound/Theorem11.lean:111}} \newline \lean{GsimOT.error_term_le}\newline{\scriptsize\texttt{Theorem11\_CutUpperBound/Theorem11.lean:58}} \newline \lean{GsimOT.dn_mul_log_two_le}\newline{\scriptsize\texttt{Theorem11\_CutUpperBound/Theorem11.lean:28}} \newline \lean{GsimOT.log_fourteen_lt_three}\newline{\scriptsize\texttt{Theorem11\_CutUpperBound/Theorem11.lean:100}} \\ \midrule
{}(not in the paper) &  & $n=1$: $\delta_\square=0$ (the paper assumes $n\ge2$ in Appendix~E.2) & \lean{GsimOT.deltaCut_le_zero_of_card_eq_one}\newline{\scriptsize\texttt{Theorem11\_CutUpperBound/Theorem11.lean:473}} \newline \lean{GsimOT.rhs_nonneg}\newline{\scriptsize\texttt{Theorem11\_CutUpperBound/Theorem11.lean:295}} \\ \midrule
{}Proof of Thm.~11 & 44 & $\varepsilon^{1-\alpha}\ln\frac e\varepsilon\ln(e\ln\frac e\varepsilon)$ is bounded on $(0,1]$, for every $\alpha<1$ & \lean{GsimOT.power_numeric}\newline{\scriptsize\texttt{Theorem11\_CutUpperBound/Theorem11.lean:200}} \newline \lean{GsimOT.power_numeric_lt_one}\newline{\scriptsize\texttt{Theorem11\_CutUpperBound/Theorem11.lean:243}} \\
\end{longtable}

\section*{Shared notions used (from \texttt{GsimOT/Common})}

\begin{longtable}{>{\raggedright\arraybackslash}p{3.2cm} >{\centering\arraybackslash}p{1.1cm} >{\raggedright\arraybackslash}p{8.4cm} >{\raggedright\arraybackslash}p{10.2cm}}
\toprule
\textbf{Paper} & \textbf{Page} & \textbf{Statement} & \textbf{Lean declaration(s) and location} \\
\midrule\endhead
\bottomrule\endfoot
{}Appendix~A & 14 & Cut norm $\|A\|_\square=\max_{S,T}|A(S,T)|$ & \lean{GsimOT.cutNorm}\newline{\scriptsize\texttt{Common/CutNorm.lean:32}} \newline \lean{GsimOT.le_cutNorm}\newline{\scriptsize\texttt{Common/CutNorm.lean:39}} \newline \lean{GsimOT.cutNorm_le}\newline{\scriptsize\texttt{Common/CutNorm.lean:44}} \newline \lean{GsimOT.exists_eq_cutNorm}\newline{\scriptsize\texttt{Common/CutNorm.lean:52}} \newline \lean{GsimOT.cutNorm_smul}\newline{\scriptsize\texttt{Common/CutNorm.lean:59}} \\ \midrule
{}Section~1, Section~2 & 2, 5 & $B(G,H,Q)$ and $C_\square(Q)=\|B(G,H,Q)\|_\square$ & \lean{GsimOT.B}\newline{\scriptsize\texttt{Common/OTDistances.lean:58}} \newline \lean{GsimOT.costCut}\newline{\scriptsize\texttt{Common/CutNorm.lean:79}} \\ \midrule
{}Section~1, Section~2 & 1--2, 5 & $d_\square$, $d_\square^{\mathrm{OT}}$ (infimum) and $\delta_\square=n^{-2}d_\square$, $\delta^{\mathrm{OT}}_\square=n^{-2}d^{\mathrm{OT}}_\square$ & \lean{GsimOT.dCut}\newline{\scriptsize\texttt{Common/CutNorm.lean:82}} \newline \lean{GsimOT.dOTCut}\newline{\scriptsize\texttt{Common/CutNorm.lean:85}} \newline \lean{GsimOT.deltaCut}\newline{\scriptsize\texttt{Common/CutNorm.lean:88}} \newline \lean{GsimOT.deltaOTCut}\newline{\scriptsize\texttt{Common/CutNorm.lean:91}} \\ \midrule
{}Appendix~E.2, Eq.~(21), E.2.2 & 32, 36 & $\zeta_R(A,B)=\|D_{A,B}\|_{\square,\mu_R}$ with $\mu_R=R/n$; $\zeta_R(A_G,A_H)=\zeta_R=n^{-2}C_\square(R)$ for graphs & \lean{GsimOT.epsR}\newline{\scriptsize\texttt{Common/CutNorm.lean:105}} \newline \lean{GsimOT.adjMat}\newline{\scriptsize\texttt{Common/CutNorm.lean:94}} \newline \lean{GsimOT.epsR_adjMat}\newline{\scriptsize\texttt{Common/CutNorm.lean:110}} \\ \midrule
{}Section~1, Appendix~B & 1, 3, 15 & $\kappa_{p\to q}=1/q+1-1/p$, $A^\pi_G-A_H$, $d_{p\to q}$ and $\delta_{p\to q}=n^{-\kappa_{p\to q}}d_{p\to q}$ ($s(n)=n^{\kappa_{p\to q}}$) & \lean{GsimOT.kappa}\newline{\scriptsize\texttt{Common/AlignedDistances.lean:23}} \newline \lean{GsimOT.alignedDiff}\newline{\scriptsize\texttt{Common/AlignedDistances.lean:31}} \newline \lean{GsimOT.distPQ}\newline{\scriptsize\texttt{Common/AlignedDistances.lean:35}} \newline \lean{GsimOT.deltaPQ}\newline{\scriptsize\texttt{Common/AlignedDistances.lean:39}} \\
\end{longtable}

\section*{Notes on the formalisation}

\begin{enumerate}[leftmargin=*]
\item \textbf{Cited results.} Theorem~35 ([HOG26, Thm.~1.4], Haqi--Oveis Gharan; p.~32) is the only remaining assumption. It is the explicit hypothesis \lean{ThinPermutationTheorem} of \lean{theorem11_of_thin}; no \texttt{axiom} is added. Its statement matches [HOG26, Thm.~1.4] (bipartite case: an $O(\log n\log\log n)$-thin perfect matching, not necessarily in $\mathrm{supp}(x)$), written with an absolute constant as $C\ln(2n)\ln\ln(4n)$ as in the paper, with thinness in the cut form of Eq.~(23). The Lean hypothesis gives thinness with exactly $\alpha=C_{\mathrm{thin}}\ln(2n)\ln\ln(4n)$, so the step ``$\alpha'$-thin with $\alpha'\le\alpha$, hence $\alpha$-thin'' of the proof of Theorem~38 is built in. The contraction principle [LT91, Thm.~4.12] (Lemma~43), also cited in the paper, is \emph{proved} (\lean{contractionPrinciple}) via Talagrand's comparison argument. The paper's other citations used here are also proved: [vdVW96, L.~2.3.1] (Lemma~42) and Carath\'eodory's theorem [Sch86] with $N\le n^2$ (proof of Lemma~48).
\item \textbf{Theorem~35 from [HOG26].} [HOG26] is formalised in the separate library \texttt{ThinMatching} (see \texttt{ThinMatching/HOG26\_correspondence.pdf}), assuming the tree-cut sparsifier theorem of R\"acke--Shah [RS14] that [HOG26] cites. \lean{thinPermutationTheorem_of_rackeShah} derives \lean{ThinPermutationTheorem} from it, so \lean{theorem11_of_rackeShah} proves Theorem~11 assuming only [RS14].
\item \textbf{Additive bound (Appendix~E.2.2): what follows the paper.} Lemma~41 is proved with the paper's support-function argument (\lean{lemma41_step}: $m_T(Z,Z)\le m_T(Z',Z)+M\eta\tau$ via $h_Z(Tu')$, then $m_T(Z',Z)\le m_T(Z',Z')+M\eta\tau$ via $h_{Z'}(T^*u)$ with the adjoint; then $T$ and $-T$); the statement does not assume $Z,Z'$ convex. Lemma~44 averages over independent uniform representatives $I_j\in B_j$ (\lean{sum_piFinset_comp}), uses Jensen, Lemma~43 for fixed $I$ and Eq.~(29); the statement does not assume the blocks disjoint. Lemma~42 is proved for arbitrary $\theta_j$. Eq.~(36) is proved as an identity (\lean{eq36}) and Corollary~46 for all $P,Q\in\mathrm{DS}(n)$. Expectations are finite weighted sums (\lean{radAvg}, \lean{bern}, uniform averages over \lean{Fintype.piFinset}).
\item \textbf{Additive bound: choices of the formalisation.} (1) The standard inner product on $\mathbb R^n$ is used instead of $\mathcal L^2(n)$, with the scalings adjusted (eigenvalues $\lambda_r/n$, eigenfunctions $\sqrt n\,e_r$); $\|K\|_{\mathcal L^2(n)}\le t$ is \lean{OpBound} $K$ $(tn)$ in the equivalent AM--GM form. (2) $H=\mathbb R^{V\sqcup W}$ (dimension $2n$) instead of $\mathbb R^{I_A\sqcup I_B}$; the extra coordinates of $z_{ij}$ are zero, so the case $I_A=I_B=\emptyset$ needs no separate treatment. (3) The blocks of Lemma~47 come from a maximal set of ``cut points'', no two of them $1$ to $L-1$ steps apart along $\tau$, instead of the explicit cycle decomposition; they satisfy the paper's properties (i) and (ii), which is all the proof uses. Lemma~44 is applied to all blocks with the vectors $X_au_i$ (zero on blocks with $X_a=0$) instead of to the blocks with $X_a=1$; the bound is the same. The w.l.o.g.\ step $P_0=I$ is done by working with $\tau=\pi_0^{-1}\pi_1$ and $u_i=z_{i,\pi_0(i)}$ directly. The proof gives $18Mn^{-1/3}$ (\lean{lemma47_sharp}), the statement $20Mn^{-1/3}$ (\lean{lemma47}). (4) In Lemma~48 the binary tree splits the index set in halves at every node (all weights positive) instead of padding with leaves of weight zero; the induction is the paper's.
\item \textbf{Theorem~38 (Appendix~E.2.1).} The proof follows the paper step by step: $\zeta_P$ as a cut norm (\lean{epsR_permMat_le}), the split $A-PBP^\top=(A-RAR^\top)+P(Q^\top AQ-B)P^\top$ with $R=PQ^\top$, the two bounds of Eq.~(27) (\lean{eq27_left}, \lean{eq27_right}), the bound on $EA$ via Lemmas~36 and~37 (\lean{normEA_le}), and Eq.~(28) with Eq.~(25) (\lean{eq28_cut}, \lean{cut_R_EAt_le}). Matrices are written entrywise ($R(v,k)=Q(k,\pi v)$). Lemma~37 is proved for $y=\mathbf1_T$, the only case used. Eq.~(25) is used in the one-dimensional form \lean{abs_sum_mul_le_of_cut}. The bound for every $\alpha$-thin $P$ (Eq.~(26)) is \lean{eq26}, the existence statement is \lean{theorem38}.
\item \textbf{Proof of Theorem~11 (Appendix~E.2.3).} The optimal $Q$ is obtained by compactness, as in the paper's ``let $Q$ be optimal''. Eq.~(39) uses $C_1=1+4C_{\mathrm{thin}}$ and $\ln(2n)\ln\ln(4n)\ge\ln4\cdot\ln\ln8>1$ for $n\ge2$ as in the paper. The case $n=1$, excluded in Appendix~E.2, is handled separately ($\delta_\square=0$). For $n>\varepsilon^{-12}$, the paper's monotonicity argument for $x\mapsto x^{-1/6}\sqrt{\ln x}$ is replaced by $\ln m\le m/e$ with $m=n^{1/6}$; the constant $112$ is the same. The constants are $C=\max\{2,39C_1,113\}$ as in the paper.
\item \textbf{Power bound.} Proved for every $\alpha<1$ (\lean{theorem11_power_lt_one_of_thin}), as stated in E.2.3; for $\alpha\le\frac12$ the bound for $\alpha=\frac12$ is used together with $f^{1/2}\le f^\alpha$ on $[0,1]$. The cut part of Theorem~2, $\delta_\square\le C_\varepsilon(\delta^{\mathrm{OT}}_\square)^{1-\varepsilon}$ for every $\varepsilon>0$, is \lean{theorem2_cut_of_thin}. (In Lean $0^\alpha=0$ for $\alpha<0$; the statement still holds since $\delta^{\mathrm{OT}}_\square=0$ forces $\delta_\square=0$.)
\item In Lean $\ln0=0$ and $e/0=0$, so the right-hand side of Theorem~11 is $0$ when $\delta^{\mathrm{OT}}_\square=0$, matching the paper's convention.
\end{enumerate}

\end{document}